% Macros and settings used by both the pdf and the web version.

% Theorem-like environments collected by the dependency graph.  They carry no
% name and no number of their own: every statement is named by the \paper or
% \aux command that precedes it, so that a statement of arXiv:2607.19651 is
% called here exactly what the paper calls it.  The environment then decides
% only the style, and which shape the dependency graph draws.
\newtheorem{theorem}{}
\newtheorem{proposition}[theorem]{}
\newtheorem{lemma}[theorem]{}
\newtheorem{corollary}[theorem]{}

\theoremstyle{definition}
\newtheorem{definition}[theorem]{}

\theoremstyle{remark}
\newtheorem{remark}[theorem]{}

% \paper{Proposition 9} names the next statement as the paper names it; the
% name is what \ref prints, so a reference reads as the paper's own.  Several
% nodes may carry one name: Remark 5 makes four claims, and each is a node.
\newcommand{\paper}[1]{\renewcommand{\thetheorem}{#1}}

% \aux{Lemma} names a statement the paper does not make --- a step it leaves
% implicit, plumbing for the sample space, a witness.  Those are numbered A.1,
% A.2, ... in their own sequence, so that no number here is ever mistaken for
% one of the paper's.
\newcounter{auxthm}
\newcommand{\aux}[1]{\stepcounter{auxthm}%
  \renewcommand{\thetheorem}{#1~A.\arabic{auxthm}}}

% Notation of the paper.
\newcommand{\actors}{\mathcal{A}}
\newcommand{\opinions}{\mathcal{O}}
\newcommand{\express}[2]{\pi^{#1,#2}}
\newcommand{\expressa}[2]{\pi_\alpha^{#1,#2}}
\newcommand{\Sspace}{S}
\newcommand{\Salpha}{S^\alpha}
\newcommand{\ladder}{L}
\newcommand{\steep}{\hat{L}}
\newcommand{\Uproc}[1]{U^{\beta,#1}}
\newcommand{\Uskel}[1]{\tilde{U}^{\beta,#1}}
\newcommand{\Rhit}[1]{R^{\beta,#1}}
\newcommand{\Rskel}[1]{\tilde{R}^{\beta,#1}}
\newcommand{\muinv}{\mu^\beta}
\newcommand{\muskel}{\tilde{\mu}^\beta}
\newcommand{\jumprate}{q_\beta}
\newcommand{\pressone}{\tfrac{1}{M-1}}
